\chapter{Seguimiento de proyecto fin de máster}
\label{Seguimiento de proyecto fin de máster}

Seguimiento del trabajo real, por oposición a la simulación de gestión empresarial que recoge el
capítulo \ref{Gestión de proyecto software}.

\section{Forma de seguimiento}

Mecanismo de seguimiento empleado: reuniones con el tutor, documentación continua del estado del
proyecto en el propio repositorio y registro de decisiones con su fecha y su justificación.

\section{Planificación inicial}

Planificación prevista al comienzo del trabajo, con sus hitos.

\section{Planificación final}

Planificación realmente ejecutada y análisis de las desviaciones. Se documentan aquí las dos
decisiones que más alteraron el plan: la reducción del conjunto de datos previsto una vez se
comprobó que el tamaño no era el factor limitante, y la incorporación del fondo fotorrealista, que
inicialmente estaba fuera de alcance.

\begin{table*}[ht]
	\centering
	\caption{Desviaciones entre la planificación inicial y la final}
	\label{tab:desviaciones}
	\makebox[\textwidth]{
		\begin{tabular}{|p{5cm}|p{5cm}|p{5cm}|}
			\cline{1-3}
			Previsto & Realizado & Causa de la desviación \\ \hline
			[PENDIENTE] & & \\ \hline
		\end{tabular}
	}
\end{table*}
